% fig2_trap_design_space — schematic: composition-tunable charge-trap design
% space beyond canonical dielectrics. figure-agent figure #2 (E2E dogfood,
% 2026-06-20). SCHEMATIC ONLY — Panel c icons are number-free qualitative shapes.
\documentclass[border=4pt]{standalone}
\usepackage{polymer-paper-preamble}
\usetikzlibrary{decorations.pathmorphing}

\begin{document}
\resizebox{178mm}{!}{%
\begin{tikzpicture}[
  every node/.style={font=\sffamily\fontsize{7}{8.4}\selectfont},
  axisArr/.style={-{Stealth[length=4pt,width=3pt]}, cGray!70!black, line width=0.7pt},
  thinArr/.style={-{Stealth[length=3.5pt,width=2.5pt]}, cGray!70!black, line width=0.6pt},
  labelStrong/.style={font=\sffamily\bfseries\fontsize{7}{8.4}\selectfont},
  labelStd/.style={font=\sffamily\fontsize{7}{8.4}\selectfont},
  labelMute/.style={font=\sffamily\itshape\fontsize{7}{8.4}\selectfont, text=cGray!80!black},
  panelLetter/.style={font=\sffamily\bfseries\fontsize{8}{9.6}\selectfont, text=cGray!90!black},
  sulfurNode/.style={circle, fill=cAmber, draw=cRed!70!black, line width=0.6pt,
                     inner sep=0pt, minimum size=3.0mm},
  canonNode/.style={circle, fill=cBlue!65, draw=cBlue!55!black, line width=0.5pt,
                    inner sep=0pt, minimum size=2.3mm},
  chargeP/.style={circle, fill=cRed!80, draw=cRed!60!black, line width=0.4pt,
                  inner sep=0pt, minimum size=1.7mm},
]

% ===================== PANEL A — origin of deep traps =====================
% Panel A
\node[panelLetter] at (0.25, 9.55) {a};
\node[labelStrong, anchor=west] at (0.55, 9.55) {Origin of deep traps};

% (1) S--S backbone with a homolytic-cleavage break (two radicals)
\draw[line width=0.9pt, cGray!85!black]
  (0.55,8.72) -- (1.05,9.07) -- (1.55,8.72) -- (2.05,9.07) -- (2.55,8.72) -- (3.02,9.05);
\draw[line width=0.9pt, cGray!85!black]
  (3.66,8.67) -- (4.05,9.07) -- (4.55,8.72);
\node[chargeP] at (3.02,9.05) {};
\node[chargeP] at (3.66,8.67) {};
\node[labelStd, text=cRed!70!black, anchor=west] at (3.18,9.30) {S$^{\bullet}$ radicals};
\node[labelMute, anchor=west] at (0.55,8.40) {S--S backbone};
\node[labelMute, anchor=north west, text width=4.3cm] at (0.55,8.08)
  {homolytic S--S bond cleavage yields a deep Coulomb trap};

% (2) "creates" arrow down to the deep well
\draw[-{Stealth[length=5pt,width=3.5pt]}, cGray!60!black, line width=0.9pt]
  (3.40,7.45) -- (3.40,6.45);
\node[labelMute, anchor=west] at (3.50,6.98) {creates};

% (3) well comparison at a shared rim baseline — direct depth contrast
\draw[line width=0.5pt, cGray!45, dash pattern=on 1.5pt off 1.5pt] (0.45,6.15) -- (4.78,6.15);
% conventional: shallow narrow well (left)
\draw[line width=0.9pt, cBlue!70!black]
  (0.55,6.15) .. controls (0.95,6.15) and (1.05,5.28) .. (1.35,5.28)
  .. controls (1.65,5.28) and (1.75,6.15) .. (2.15,6.15);
\node[chargeP, minimum size=1.4mm, fill=cBlue!70, draw=cBlue!55!black] at (1.35,5.45) {};
\node[labelStrong, text=cBlue!60!black, anchor=north, align=center, text width=2.0cm] at (1.35,5.00)
  {conventional\\high $\varepsilon$};
\node[labelMute, anchor=north, align=center, text width=2.0cm] at (1.35,4.40)
  {shallow well, charge detraps};
% sulfur: wide deep well (right)
\draw[line width=1.1pt, cRed!70!black]
  (2.55,6.15) .. controls (3.05,6.15) and (3.05,2.60) .. (3.66,2.60)
  .. controls (4.27,2.60) and (4.27,6.15) .. (4.72,6.15);
\node[chargeP] at (3.66,2.85) {};
\node[chargeP] at (3.41,2.97) {};
\node[chargeP] at (3.91,2.97) {};
\node[labelStrong, text=cRed!70!black, anchor=north, align=center, text width=2.3cm] at (3.64,2.35)
  {sulfur polymer\\low $\varepsilon$};
\node[labelMute, anchor=north, align=center, text width=2.4cm] at (3.64,1.78)
  {wide, deep well, charge retained};

% panel divider
\draw[line width=0.4pt, cGray!35] (4.95,0.45) -- (4.95,9.75);

% ===================== PANEL B — trap design space (HERO) =================
% Panel B
\node[panelLetter] at (5.25, 9.55) {b};
\node[labelStrong, anchor=west] at (5.55, 9.55) {Composition-tunable charge-trap design space};

% axes
\draw[axisArr] (6.10,3.90) -- (17.25,3.90);
\draw[axisArr] (6.10,3.90) -- (6.10,9.15);
\node[labelStd, anchor=north] at (11.70,3.62) {trap-distribution breadth\quad($n$)};
\node[labelStd, anchor=south, rotate=90] at (5.78,6.50) {charge retention \& resistivity\quad($\rho$)};

% beyond-conventional region
\fill[cAmber!12, rounded corners=3mm] (10.40,5.70) rectangle (17.00,8.95);
\node[labelMute, text=cAmber!55!black, anchor=north east, text width=3.6cm, align=right]
  at (16.90,8.85) {beyond conventional dielectrics};

% conventional cluster
\fill[cBlue!10, rounded corners=2mm] (6.45,4.15) rectangle (8.75,5.45);
\node[canonNode] at (7.00,4.70) {};
\node[canonNode] at (7.75,4.55) {};
\node[canonNode] at (7.45,5.05) {};
\node[labelMute, anchor=south, text width=2.6cm, align=center] at (7.60,5.58)
  {PI, PDMS, PET (shallow, leaky)};

% sulfur S60->S85 trajectory (qualitative arc — NOT measured coordinates)
\draw[line width=1.4pt, cAmber!85!black, -{Stealth[length=7pt,width=5pt]}]
  plot[smooth, tension=0.7]
  coordinates {(9.50,5.40)(10.90,6.75)(12.40,7.55)(14.00,7.65)(15.60,6.95)};
\foreach \px/\py in {9.50/5.40, 10.90/6.75, 12.40/7.55, 14.00/7.65, 15.60/6.95}{%
  \node[sulfurNode] at (\px,\py) {};}
\node[labelStd, anchor=north east] at (9.40,5.25) {S60};
\node[labelStd, anchor=west] at (15.80,6.95) {S85};
\node[labelStrong, text=cRed!65!black, anchor=south] at (12.60,7.85) {sulfur content $\uparrow$};

% ===================== PANEL C — kinetic signature (schematic) ============
% Panel C
\node[panelLetter] at (5.25, 3.15) {c};
\node[labelStrong, anchor=west] at (5.55, 3.15) {Kinetic charge-trapping signature (schematic)};

% left icon: I(t) ~ t^-n  (log-log slope = -n; sulfur high n = steeper)
\draw[thinArr] (6.55,0.72) -- (6.55,2.82);
\draw[thinArr] (6.55,0.72) -- (10.45,0.72);
\node[labelMute, anchor=south, rotate=90] at (6.30,1.80) {log $I(t)$};
\node[labelMute, anchor=north] at (8.45,0.64) {log $t$};
% conventional: shallow slope (low n)
\draw[line width=1.0pt, cBlue!70!black] (6.85,2.55) -- (10.10,2.08);
% sulfur: steep slope (high n)
\draw[line width=1.3pt, cRed!75!black] (6.85,2.55) -- (9.05,0.98);
% Debye reference (exponential — bends down fastest at long t)
\draw[line width=0.7pt, cGray!60!black, dash pattern=on 2.5pt off 1.5pt]
  (6.85,2.55) .. controls (7.55,2.30) and (7.75,1.35) .. (8.05,0.95);
\node[labelStd, text=cBlue!65!black, anchor=west] at (8.55,2.60) {conventional (low $n$)};
\node[labelStd, text=cRed!70!black, anchor=west] at (8.85,1.55) {sulfur (high $n$)};
\node[labelMute, anchor=east] at (7.90,1.05) {Debye};

% right icon: rho vs sulfur content (magnitude, independent metric)
\draw[thinArr] (11.95,0.72) -- (11.95,2.82);
\draw[thinArr] (11.95,0.72) -- (16.95,0.72);
\node[labelMute, anchor=south, rotate=90] at (11.70,1.80) {$\rho$};
\node[labelMute, anchor=north] at (14.45,0.64) {sulfur content};
\fill[cBlue!18, rounded corners=1mm] (12.20,0.90) rectangle (16.55,1.08);
\node[labelMute, text=cBlue!60!black, anchor=west] at (12.30,0.99) {conventional dielectrics};
\draw[line width=1.2pt, cAmber!85!black]
  plot[smooth, tension=0.6]
  coordinates {(12.40,1.40)(13.40,2.00)(14.40,2.40)(15.40,2.48)(16.35,2.28)};
\foreach \qx/\qy in {12.40/1.40, 13.40/2.00, 14.40/2.40, 15.40/2.48, 16.35/2.28}{%
  \node[sulfurNode, minimum size=2.2mm] at (\qx,\qy) {};}

\end{tikzpicture}%
}
\end{document}
